\section{基础模型与参数来源}
\label{cr:model}
\subsection{几何、电路与物理假设}
基础模型为自编二维平面NMOS加理想存储电容的1T1C单元。设栅长为$L$，
硅体宽度为$L+0.6\um$、深度为$1.2\um$，二维有效第三维厚度为$1\um$。
主模型氧化层厚度为$5\,\mathrm{nm}$，体区为P型，体掺杂浓度
$10^{17}\,\mathrm{cm}^{-3}$，源漏高斯峰值为$10^{20}\,\mathrm{cm}^{-3}$，
结深参数为$0.12\um$。栅功函数设为$4.6\,\mathrm{eV}$，
温度为$300\,\mathrm{K}$，外接存储电容为$10\,\mathrm{fF}$。
比较的四个栅长是$0.18/0.25/0.35/0.45\um$，表示访问晶体管的几何长度，
并不表示商业DRAM工艺节点。

模型使用Fermi统计、OldSlotboom带隙修正、掺杂和电场相关迁移率以及
SRH、Auger复合。器件采用漂移扩散求解，电路节点同时参与瞬态计算。
选用这种结构的依据是它保留了访问管、电荷存储及电荷共享三项基本机制，
输入规模适合当前计算环境，也便于检查粒子源与电路响应之间的关系。
真实DRAM可能采用复杂沟道、埋栅和立体电容，简化结构不能替代这些工艺的
电学标定，更不能由栅长数值推断产品代际。

\begin{table}[htbp]
\centering\small
\caption{主模型参数与证据性质。数值为本项目冻结的探索参数，不能统一称为文献实测值。}
\label{cr:parameters}
\begin{tabularx}{\textwidth}{lllX}
\toprule
参数&主模型取值&用途&来源或适用范围\\
\midrule
栅长$L$&0.18至0.45 $\mu$m&几何扫描&自编平面访问管，四点离散取值\\
氧化层&5 nm&主矩阵&4/6 nm为后续受控变体，尚无最终结果\\
存储电容&10 fF&储存电荷&1T1C原理与已有混合模式示例支持形式，未校准器件\\
位线与参考电容&100 fF&读出电荷共享&统一电路假设，便于匹配比较\\
初态SN&0.1/0.9 V&存储0/1&理想钳位初始化，非真实写入校准\\
中性参考SN&0.5 V&消除共模偏移&同器件、同网格与时序的参考支路\\
预充导通电阻&50 $\Omega$&R50-main读出&由独立probe支持的物理修订\\
\bottomrule
\end{tabularx}
\end{table}

\subsection{从参数来源到可检验假设}
参数来源分为原理、工具约定、已有模型形式和本项目假设四类。
DRAM教材支持访问晶体管与电容的基本工作机制；Geant4资料支持输运过程和
电磁模型的含义\cite{keeth2001,g4em}。已有混合模式DRAM示例说明外接理想
电容的建模方式可行，但本项目没有复制并校准其完整工艺工程。
几何尺寸、掺杂、位线电容及裕量门槛仍按探索假设处理。用明确的来源分层，
可以把“软件支持这个模型”与“模型代表某种真实器件”区分开来。

R50-main的0/1状态通过理想SN钳位获得，钳位在$0.6\ns$释放，
并在$0.95\ns$检查初态误差不超过$1\mV$。这一设置减少不同栅长写入程度
不同所带来的混杂，使比较更接近“相同初态下的访问管几何变化”。
然而它只控制初始电压，不证明实际写入时间、能耗或访问管写入能力。
v1曾采用实际访问管充电，不同栅长的粒子前电压不同；
因此两版数据适合解释设计改进，不能拼接成一条统一的尺寸响应曲线。

\subsection{保持、预充与匹配参考}
粒子在$1\ns$注入，保持指标在$2.8\ns$提取。位线随后经预充支路升至
$0.5\V$，在$3\ns$释放；字线在$3.05$至$3.10\ns$升至$1.8\V$，
读出在$4\ns$采样。R50-main使用$50\,\Omega$导通电阻，
关断电阻为$10^{20}\,\Omega$，不把理想开关的电压跳变当作预充完成。

匹配中性参考在相同尺寸、网格、偏压及时序下取SN=$0.5\V$。
批量计算使用信号位线与中性参考位线之差，并先通过真实双支路仿真核对。
这种参考随访问管打开而产生相应的耦合和电荷交换，比统一减去一个常数更适合
处理状态相关偏移。参考中不注入粒子；物理双支路中采用独立器件定义，
以防同一HeavyIon设置同时作用于信号与参考。
